\documentclass[11pt]{article}
\usepackage[margin=1in]{geometry}
\usepackage[T1]{fontenc}
\usepackage[utf8]{inputenc}
\usepackage{lmodern}
\usepackage{microtype}
\usepackage{booktabs}
\usepackage{enumitem}
\usepackage{hyperref}
\setlength{\parindent}{0pt}
\setlength{\parskip}{0.6em}

\title{Prompt-Profile Build Recipe}
\author{}
\date{2026-03-23 06:42 UTC}

\begin{document}
\maketitle

\section*{Bottom line}
Treat the next prompt-prefill predictor as a prompt-level terminal-statistics model,
not as another thresholded loop classifier.

\begin{itemize}[leftmargin=1.5em]
\item If exactly one head ships, ship \texttt{mean\_relative\_length = E[L / E]}.
\item The default bundle should still include \texttt{p\_loop} as the cleaner
failure-prox companion head.
\item Use the per-layer ensemble with late aggregation as the default readout.
\item Ship the ranking-oriented checkpoint (\texttt{best\_rank}), not the loss-best checkpoint.
\end{itemize}

\section*{Why this is the right current default}
The finished direct-head table already separates the ``useful'' scalar from the
``cleaner failure'' scalar:

\begin{itemize}[leftmargin=1.5em]
\item \textbf{\texttt{mean\_relative\_length}} is the strongest and steadiest head on the
finished same-archive runs:
  \begin{itemize}[leftmargin=1.5em]
  \item \texttt{GPQA}: Spearman \texttt{0.649}
  \item \texttt{AIME}: Spearman \texttt{0.769} at \texttt{best\_loss}, \texttt{0.853} at \texttt{best\_rank}
  \item \texttt{MATH-500}: Spearman \texttt{0.716}
  \item \texttt{MMLU-Pro}: Spearman \texttt{0.303}
  \end{itemize}
\item \textbf{\texttt{p\_loop}} stays worth training because it is more failure-proximal and
less prompt-length-shaped on the heavier-tail slices:
  \begin{itemize}[leftmargin=1.5em]
  \item \texttt{GPQA}: Spearman \texttt{0.371}, top-\texttt{20\%} capture \texttt{0.318}
  \item \texttt{AIME}: Spearman \texttt{0.475}, top-\texttt{20\%} capture \texttt{0.8}
  \item \texttt{MMLU-Pro}: Spearman \texttt{0.335}, top-\texttt{20\%} capture \texttt{0.652}
  \end{itemize}
\item \textbf{\texttt{loop\_budget\_share}} did not survive the lower-tail controls cleanly
enough to replace the current bundle.
\item \textbf{\texttt{majority\_s\_t}} remains control-only: on \texttt{AIME}, prompt length
alone already gets \texttt{PR-AUC 0.898} for \texttt{majority\_s\_0.5}.
\end{itemize}

That leaves a simple practical split:
\begin{itemize}[leftmargin=1.5em]
\item \texttt{mean\_relative\_length}: use when the requirement is ``be useful''.
\item \texttt{p\_loop}: keep in the bundle so the system still scores something close to the
actual failure event.
\end{itemize}

\section*{Exact next-build recipe}
\textbf{Inputs}
\begin{itemize}[leftmargin=1.5em]
\item prompt-prefill activations only;
\item one selected layer or per-layer independent heads with late aggregation;
\item do \textbf{not} go back to stacking all layers into one vector as the main surface.
\end{itemize}

\textbf{Decode policy}
\begin{itemize}[leftmargin=1.5em]
\item fixed \texttt{temperature = 0.2};
\item fixed \texttt{num\_generations};
\item fixed \texttt{max\_tokens} and \texttt{max\_model\_len};
\item prompt-disjoint, in-distribution train/test split.
\end{itemize}

\textbf{Targets}
\begin{enumerate}[leftmargin=1.5em]
\item Build one repeated-rollout archive per prompt.
\item Relabel it to:
  \begin{itemize}[leftmargin=1.5em]
  \item \texttt{--target-kind regression --profile-target mean\_relative\_length}
  \item \texttt{--target-kind probability --profile-target p\_loop}
  \end{itemize}
\item Train the ensemble view by default and keep last-layer as the cheap control.
\end{enumerate}

\textbf{Checkpointing}
\begin{itemize}[leftmargin=1.5em]
\item keep \texttt{best\_loss} for calibration-style reporting;
\item ship \texttt{best\_rank} for utility-facing prompt ranking.
\end{itemize}

\section*{Required baselines and success criteria}
Every run now needs three non-activation baselines:
\begin{itemize}[leftmargin=1.5em]
\item \texttt{prompt\_length} only;
\item \texttt{effective\_budget} only;
\item joint \texttt{prompt\_length + effective\_budget}.
\end{itemize}

The acceptance rule is not ``beat one scalar baseline on one metric.''

\begin{itemize}[leftmargin=1.5em]
\item For \texttt{mean\_relative\_length}, usefulness means:
  \begin{itemize}[leftmargin=1.5em]
  \item primary metric: held-out prompt Spearman;
  \item secondary metric: top-\texttt{20\%} capture of the bad tail;
  \item guardrail: MSE.
  \end{itemize}
\item For \texttt{p\_loop}, usefulness means:
  \begin{itemize}[leftmargin=1.5em]
  \item primary metric: top-\texttt{20\%} capture of looping prompts;
  \item secondary metric: Spearman against empirical \texttt{p\_loop};
  \item guardrail: Brier.
  \end{itemize}
\end{itemize}

Do not call a head useful unless it beats the strongest non-degenerate baseline on
its primary ranking metric.

\section*{What to de-prioritize now}
\begin{itemize}[leftmargin=1.5em]
\item \texttt{majority\_s\_t} as a shipped objective: too leakage-shaped.
\item \texttt{loop\_budget\_share} as the default scalar: too brittle outside the heavy tail.
\item arbitrary threshold hunting over \texttt{s\_t}: it throws away information and
invites length leakage.
\item stacked-all-layer features as the main readout surface.
\item loss-only checkpoint selection as the default ship criterion.
\end{itemize}

\section*{One fallback if LiveCodeBench contradicts the current ordering}
If the unfinished \texttt{LiveCodeBench} lane comes back badly enough that one more
binary head needs to be reopened, test direct \texttt{p\_cap} next.

Do \textbf{not} reopen another thresholded \texttt{s\_t} sweep or promote
\texttt{loop\_budget\_share} ahead of that.

\end{document}
